% GENERATED FILE. Edit canonical lessons, then run npm run generate:books.
% canonical-source-hash: fnv1a64:99167616db240d40
% canonical-lessons: GE-C514-melodie, GE-C514-rhythmus, GE-C514-takt, GE-C514-strophe, GE-C514-summen

\chapter{Melodie und Rhythmus}
\label{ch:ge-melodie-und-rhythmus}
\hlchaptermodality{514}

\begin{chapteropening}
\textbf{By the end of this chapter:} \emph{I can talk about tunes, rhythm and songs.}

Complete the last lesson of chapter 514: I can talk about tunes, rhythm and songs.
\end{chapteropening}

\section[die Melodie]{die Melodie — melody, tune}

\label{lesson:GE-C514-melodie}



\begin{quote}
Before the new one: say the German for drums, then the German for to drum.
\end{quote}

\subsection*{You'll want to know: die Melodie}
\textbf{die Melodie} — \textquotedblleft{}melody, tune\textquotedblright{}.

\begin{quote}
Die Melodie von diesem Lied ist sehr schön.
\end{quote}

\begin{quote}
\emph{The tune of this song is very beautiful.}
\end{quote}

Plural: \emph{die Melodien}.

\begin{cousinweb}
\emph{Melodie} was borrowed, through Late Latin \emph{melōdia}, from Greek \emph{melōidía} \textquotedblleft{}singing, a song\textquotedblright{}; EtymDB lists the German word without a source.
\end{cousinweb}

\begin{grammarlens}[title={What you've built}]
One more word for telling, explaining and describing at B1.
\end{grammarlens}

\subsection*{Your turn}
\begin{itemize}
\raggedright
  \item \emph{Say it:} \emph{die Melodie}
  \item \emph{Say it:} \emph{die Melodie}, once more
  \item \emph{Say it:} a sentence of your own with \emph{die Melodie}
  \item \emph{Recall:} say the German for drums, then the German for to drum, then say \emph{die Melodie} again
\end{itemize}

\begin{tcolorbox}[breakable,colback=teal!4,colframe=teal!35!black,title={Before you move on}]
What does die Melodie mean? (\textquotedblleft{}Melody, tune\textquotedblright{}.) Say it once more.
\end{tcolorbox}
\section[der Rhythmus]{der Rhythmus — rhythm}

\label{lesson:GE-C514-rhythmus}



\begin{quote}
Before the new one: say the German for to drum, then the German for melody.
\end{quote}

\subsection*{You'll want to know: der Rhythmus}
\textbf{der Rhythmus} — \textquotedblleft{}rhythm\textquotedblright{}.

\begin{quote}
Die Kinder tanzen im Rhythmus der Musik.
\end{quote}

\begin{quote}
\emph{The children dance to the rhythm of the music.}
\end{quote}

Plural: \emph{die Rhythmen}.

\begin{cousinweb}
\emph{Rhythmus} is a learned borrowing from Latin \emph{rhythmus}, from Greek \emph{rhythmós} \textquotedblleft{}measured flow, time\textquotedblright{} (Lewis \& Short).
\end{cousinweb}

\begin{grammarlens}[title={What you've built}]
One more word for telling, explaining and describing at B1.
\end{grammarlens}

\subsection*{Your turn}
\begin{itemize}
\raggedright
  \item \emph{Say it:} \emph{der Rhythmus}
  \item \emph{Say it:} \emph{der Rhythmus}, once more
  \item \emph{Say it:} a sentence of your own with \emph{der Rhythmus}
  \item \emph{Recall:} say the German for to drum, then the German for melody, then say \emph{der Rhythmus} again
\end{itemize}

\begin{tcolorbox}[breakable,colback=teal!4,colframe=teal!35!black,title={Before you move on}]
What does der Rhythmus mean? (\textquotedblleft{}Rhythm\textquotedblright{}.) Say it once more.
\end{tcolorbox}
\section[der Takt]{der Takt — beat, time (in music); tact}

\label{lesson:GE-C514-takt}



\begin{quote}
Before the new one: say the German for melody, then the German for rhythm.
\end{quote}

\subsection*{You'll want to know: der Takt}
\textbf{der Takt} — \textquotedblleft{}beat, time (in music); tact\textquotedblright{}.

\begin{quote}
Der Schlagzeuger hält den Takt.
\end{quote}

\begin{quote}
\emph{The drummer keeps the beat.}
\end{quote}

Plural: \emph{die Takte}. Also tact: \emph{Er hat keinen Takt.}

\begin{cousinweb}
\emph{Takt} comes from Latin \emph{tāctus} \textquotedblleft{}touch\textquotedblright{} (EtymDB), from \emph{tangere} \textquotedblleft{}to touch\textquotedblright{}.
\end{cousinweb}

\begin{grammarlens}[title={What you've built}]
One more word for telling, explaining and describing at B1.
\end{grammarlens}

\subsection*{Your turn}
\begin{itemize}
\raggedright
  \item \emph{Say it:} \emph{der Takt}
  \item \emph{Say it:} \emph{der Takt}, once more
  \item \emph{Say it:} a sentence of your own with \emph{der Takt}
  \item \emph{Recall:} say the German for melody, then the German for rhythm, then say \emph{der Takt} again
\end{itemize}

\begin{tcolorbox}[breakable,colback=teal!4,colframe=teal!35!black,title={Before you move on}]
What does der Takt mean? (\textquotedblleft{}Beat, time (in music); tact\textquotedblright{}.) Say it once more.
\end{tcolorbox}
\section[die Strophe]{die Strophe — verse, stanza}

\label{lesson:GE-C514-strophe}



\begin{quote}
Before the new one: say the German for rhythm, then the German for beat.
\end{quote}

\subsection*{You'll want to know: die Strophe}
\textbf{die Strophe} — \textquotedblleft{}verse, stanza\textquotedblright{}.

\begin{quote}
Wir singen nur die erste Strophe.
\end{quote}

\begin{quote}
\emph{We only sing the first verse.}
\end{quote}

Plural: \emph{die Strophen}. The chorus is \emph{der Refrain}.

\begin{cousinweb}
\emph{Strophe} comes from Latin \emph{stropha}, from Greek \emph{stroph\'{\={e}}} \textquotedblleft{}a turning\textquotedblright{} (EtymDB).
\end{cousinweb}

\begin{grammarlens}[title={What you've built}]
One more word for telling, explaining and describing at B1.
\end{grammarlens}

\subsection*{Your turn}
\begin{itemize}
\raggedright
  \item \emph{Say it:} \emph{die Strophe}
  \item \emph{Say it:} \emph{die Strophe}, once more
  \item \emph{Say it:} a sentence of your own with \emph{die Strophe}
  \item \emph{Recall:} say the German for rhythm, then the German for beat, then say \emph{die Strophe} again
\end{itemize}

\begin{tcolorbox}[breakable,colback=teal!4,colframe=teal!35!black,title={Before you move on}]
What does die Strophe mean? (\textquotedblleft{}Verse, stanza\textquotedblright{}.) Say it once more.
\end{tcolorbox}
\section[summen]{summen — to hum; to buzz}

\label{lesson:GE-C514-summen}



\begin{quote}
Before the new one: say the German for beat, then the German for verse.
\end{quote}

\subsection*{You'll want to know: summen}
\textbf{summen} — \textquotedblleft{}to hum; to buzz\textquotedblright{}.

\begin{quote}
Beim Kochen summt meine Mutter ein Lied.
\end{quote}

\begin{quote}
\emph{My mother hums a song while she cooks.}
\end{quote}

Also of bees: \emph{Die Bienen summen.}

\begin{cousinweb}
\emph{summen} is inherited from Middle High German \emph{summen} (EtymDB), an imitative word for a humming sound.
\end{cousinweb}

\begin{grammarlens}[title={What you've built}]
One more word for telling, explaining and describing at B1.
\end{grammarlens}

\subsection*{Your turn}
\begin{itemize}
\raggedright
  \item \emph{Say it:} \emph{summen}
  \item \emph{Say it:} \emph{summen}, once more
  \item \emph{Say it:} a sentence of your own with \emph{summen}
  \item \emph{Recall:} say the German for beat, then the German for verse, then say \emph{summen} again
\end{itemize}

\begin{tcolorbox}[breakable,colback=teal!4,colframe=teal!35!black,title={Before you move on}]
What does summen mean? (\textquotedblleft{}To hum; to buzz\textquotedblright{}.) Say it once more.
\end{tcolorbox}
